\chapter{Adaptive Mesh Refinement}
\label{ch:amr}

\chapterorient{\Cref{ch:lifecycle} showed the simulation lifecycle closing into a
loop: solve, estimate the error, mark the worst elements, refine the mesh, solve
again. We called that the adaptive estimate--mark--refine loop and named it as our
first \emph{fold}, but we left its three working stages as black boxes. This
chapter opens them. We ask three questions in turn. \emph{Where} is the error
worst---when we do not know the true solution and so cannot subtract it? \emph{Which}
elements should we refine---given a per-element error guess, how many, and which?
And \emph{how} does the refinement happen, and why does the whole loop terminate
with a good answer using far fewer unknowns than a uniform mesh would need? Two of
the three answers are short, sharp algorithms we can run by hand---a flux-recovery
estimate and a bulk-marking rule---and the third is a mesh operation we will treat
as a boundary. By the end you will be able to compute an error indicator on a tiny
mesh, mark elements by the D\"orfler rule on a list of numbers, and read a
convergence curve that shows adaptivity earning its keep.}

\section{Why adapt at all}
\label{sec:why-adapt}

Refining a mesh means cutting its elements into smaller ones, raising the number of
unknowns $N$ and---by C\'ea's lemma (\cref{ch:galerkin})---driving the discretization
error down. The simplest way to refine is \emph{uniformly}: cut every element. In
one step a 2-D mesh of $N$ elements becomes $4N$; in three dimensions, $8N$. Uniform
refinement is easy to implement and easy to reason about, and for a perfectly smooth
solution on a simple domain it is hard to beat. But the solutions an engineering
simulator actually computes are rarely smooth everywhere.

Real fields have \emph{features}. An electrostatic potential develops a singularity
at a sharp conductor corner---the field strength blows up like $r^{-\gamma}$ as the
distance $r$ to the corner goes to zero. A magnetic field crowds into the gap between
two poles. The current concentrates at a material interface where the permittivity
jumps. Near such a feature the solution bends violently and the piecewise-polynomial
basis of \cref{ch:galerkin} struggles to follow it, so the local error is large.
Away from the feature---across the broad, gently-varying interior---the same basis is
already more than accurate enough, and the local error is tiny.

\Cref{fig:cornermesh} is the canonical picture. A re-entrant (``L-shaped'') corner
pulls a singularity into the field; a good mesh is \emph{fine} in a small
neighbourhood of the corner and \emph{coarse} everywhere else. Uniform refinement,
by contrast, would spend the same effort everywhere---lavishing resolution on the
smooth interior, where it buys almost nothing, in order to resolve the corner.

\begin{figure}[t]
\centering
\begin{tikzpicture}[scale=1.0, every node/.style={font=\footnotesize}]
  % L-shaped domain outline (re-entrant corner at the origin)
  \def\a{4.0}
  \draw[simInk, thick, fill=simBgGray]
    (0,0) -- (\a,0) -- (\a,2.0) -- (2.0,2.0) -- (2.0,\a) -- (0,\a) -- cycle;
  % the re-entrant corner
  \fill[simRust] (2.0,2.0) circle (2.2pt);
  \node[simRust, above right=0pt and 1pt] at (2.0,2.0) {re-entrant corner};
  % graded mesh: coarse far cells
  \draw[simBlue!55, thin] (0,0) grid[step=1.0] (\a,2.0);
  \draw[simBlue!55, thin] (0,2.0) grid[step=1.0] (2.0,\a);
  % progressively finer cells toward the corner: a few nested rings
  \foreach \s in {0.5,0.25,0.125}{
    \draw[simBlue!75, thin] (2.0-\s,2.0-\s) rectangle (2.0,2.0);
    \draw[simBlue!75, thin] (2.0-\s,2.0) rectangle (2.0,2.0+\s);
    \draw[simBlue!75, thin] (2.0,2.0-\s) rectangle (2.0+\s,2.0);
  }
  % subdivide those finer rings further (just visual: cross lines)
  \draw[simBlue!75, thin] (1.5,2.0)--(2.0,2.0); % already there
  % field-strength annotation
  \draw[refedge, simRust] (3.4,3.0) to[bend right=20] (2.15,2.15);
  \node[simRust, font=\scriptsize, align=center] at (3.55,3.2)
    {$|\vE|\sim r^{-\gamma}$\\(error large)};
  \node[simGray, font=\scriptsize, align=center] at (1.0,1.0)
    {smooth interior\\(error tiny)};
\end{tikzpicture}
\caption[A mesh graded to a corner singularity]{The canonical adaptive mesh. A
re-entrant corner (here the inner corner of an L-shaped domain) forces the field to
blow up like $r^{-\gamma}$ as the distance $r$ to the corner shrinks, so the
discretization error concentrates there. A well-adapted mesh is heavily refined in a
small neighbourhood of the corner---the nested rings of ever-smaller cells---and
left coarse across the smooth interior, where the basis already resolves the field.
Uniform refinement would instead cut \emph{every} cell, spending most of its new
unknowns on the interior, where they buy almost nothing.}
\label{fig:cornermesh}
\end{figure}

The payoff is captured by a single principle. C\'ea's lemma told us the total error
is governed by how well the basis can approximate the solution; interpolation theory
adds that, for a fixed total number of unknowns, the approximation error is
\emph{minimised} when the local error contributions are \emph{equal} across
elements---when the error is \emph{equidistributed}. A uniform mesh equidistributes
the cell \emph{size}, not the error: it leaves large error at the corner and
negligible error in the interior, a wildly unequal distribution. An adaptive mesh
equidistributes the \emph{error} by making cells small exactly where the error per
cell is large. The result is dramatic: for a corner singularity in two dimensions, a
uniform mesh achieves error $\mathcal{O}(N^{-\gamma/2})$ with a stalled exponent set
by the singularity strength $\gamma$, while a properly graded adaptive mesh recovers
the full smooth-solution rate $\mathcal{O}(N^{-p/2})$---the same accuracy for
\emph{orders of magnitude} fewer unknowns. \Cref{sec:convergence} makes this
quantitative; \cref{fig:convergence} draws it.

\begin{keyidea}
Uniform refinement equidistributes \emph{cell size}; adaptive refinement
equidistributes \emph{error}. Because a solution's error concentrates at features
(corners, interfaces, field concentrations) and is negligible across smooth regions,
putting small cells only where the error is large reaches a target accuracy with far
fewer unknowns than cutting every cell. The recipe to do this automatically is the
loop of \cref{ch:lifecycle}: \emph{estimate} where the error is, \emph{mark} a bulk
fraction of it, \emph{refine} those elements, and \emph{repeat}. The rest of this
chapter is those four words.
\end{keyidea}

There is a circularity to defeat first. To put small cells where the error is large,
we must know where the error is large---but the error is $u - u_h$, and we do not
know the exact solution $u$ (if we did, we would not be running the simulation). The
whole edifice rests on \emph{estimating} the error from the computed solution alone,
without ever seeing the truth. That is the subject of the next section, and it is the
cleverest of the three stages.

\section{Estimate: Zienkiewicz--Zhu flux recovery}
\label{sec:estimate}

We have a computed solution $u_h$ and want a number $\eta_K \ge 0$ for each element
$K$---an \emph{error indicator}---large where the local error $\norm{u - u_h}_K$ is
large and small where it is small. We cannot compute the error directly. The
Zienkiewicz--Zhu (ZZ) idea is to find a \emph{computable quantity that is zero
exactly when the error is zero} and grows with it---a proxy we \emph{can} evaluate.

\subsection{The discontinuous flux and the smoothing trick}

The proxy is built from the \emph{flux}. In the electrostatic problem the field is
$\vE = -\Grad V$ and the physical flux is the displacement $\vD = \varepsilon\vE$;
in the magnetostatic problem the flux is $\vH = \mu^{-1}\vB$. Generically write the
flux as $\mat F = Q\,\diffop u_h$, a material weight $Q$ times a derivative of the
discrete solution. Here is the key observation about this flux as the finite-element
method computes it.

The discrete solution $u_h$ is continuous across element boundaries (it is built from
the conforming basis of \cref{ch:galerkin}), but its \emph{derivative} is not. A
piecewise-linear $u_h$ has a constant gradient on each element that jumps as you
cross into the next. So the computed flux $\mat F = Q\,\diffop u_h$ is
\emph{discontinuous across element faces}---it has jumps. The \emph{true} flux of the
exact solution, by contrast, is smooth (continuity of the normal flux is a physical
conservation law). The size of the jumps in the computed flux therefore measures how
far the discrete solution is from the truth: where $u_h$ resolves the solution well,
its flux is nearly continuous and the jumps are small; where $u_h$ struggles, the
flux is wildly discontinuous and the jumps are large.

ZZ turns this qualitative observation into a number with one move: \emph{recover a
smooth flux} $G$ from the discontinuous one and measure the gap. Concretely, $G$ is
the $L^2$ projection of the computed flux $\mat F$ onto a conforming
(continuous-where-it-should-be) finite-element space:
\begin{equation}
G = \arg\min_{g \in V_{\text{smooth}}} \int_\dom \norm{g - \mat F}^2 \dV,
\qquad\text{equivalently}\qquad
G = \mat M^{-1}\,\bigl(\text{Flux}\cdot u_h\bigr),
\label{eq:zz-recover}
\end{equation}
where $\mat M$ is the mass matrix of the smooth space and $\text{Flux}\cdot u_h$
assembles the projection's right-hand side---a single mass-matrix solve, exactly the
kind of system the solvers of \cref{ch:krylov} dispatch. The smoothed flux $G$ is the
code's best guess at the \emph{true}, continuous flux; the gap $\mat F - G$ between
the rough computed flux and its smoothed version stands in for the unknowable error.

\begin{figure}[t]
\centering
\begin{tikzpicture}[scale=1.0]
\begin{axis}[
  width=12.5cm, height=5.4cm,
  xmin=-0.05, xmax=4.05, ymin=-0.15, ymax=2.6,
  xtick={0,1,2,3,4},
  xticklabels={$x_0$,$x_1$,$x_2$,$x_3$,$x_4$},
  ytick=\empty, ylabel={flux}, axis lines=left,
  tick label style={font=\footnotesize},
  legend style={at={(0.97,0.97)}, anchor=north east, font=\footnotesize,
    draw=simGray!50},
  clip=false,
]
  % discontinuous element flux: constant per element, jumps at nodes
  \addplot[simRust, very thick] coordinates {(0,0.35)(1,0.35)};
  \addplot[simRust, very thick] coordinates {(1,0.75)(2,0.75)};
  \addplot[simRust, very thick] coordinates {(2,1.7)(3,1.7)};
  \addplot[simRust, very thick] coordinates {(3,2.15)(4,2.15)};
  \addlegendentry{$\mathbf{F}$ (computed, discontinuous)}
  % the jumps (dotted vertical risers)
  \addplot[simRust, densely dotted, thin] coordinates {(1,0.35)(1,0.75)};
  \addplot[simRust, densely dotted, thin] coordinates {(2,0.75)(2,1.7)};
  \addplot[simRust, densely dotted, thin] coordinates {(3,1.7)(3,2.15)};
  % recovered smooth flux: a continuous (piecewise-linear through midpoints) curve
  \addplot[simBlue, very thick, mark=*, mark size=1.4pt] coordinates
    {(0,0.2)(0.5,0.35)(1.5,0.75)(2.5,1.7)(3.5,2.15)(4,2.35)};
  \addlegendentry{$G$ (recovered, smooth)}
  % gap annotation on the big-jump element (between x2 and x3)
  \draw[<->, simGreen, thick] (axis cs:2.5,1.30) -- (axis cs:2.5,1.70);
  \node[simGreen, font=\scriptsize, anchor=west] at (axis cs:2.58,1.5)
    {gap $=\eta_K$};
  \node[simGray, font=\scriptsize, anchor=north] at (axis cs:0.5,0.30)
    {small gap};
\end{axis}
\end{tikzpicture}
\caption[Zienkiewicz--Zhu flux recovery]{The ZZ estimator in one dimension. The
computed flux $\mathbf F = Q\,\mathcal{D}u_h$ (rust) is \emph{piecewise constant} and
\emph{jumps} across each element boundary, because $u_h$'s derivative is
discontinuous. The recovered flux $G$ (blue) is the smooth $L^2$ projection of
$\mathbf F$ onto a conforming space---the code's best estimate of the true,
continuous flux. The per-element indicator $\eta_K$ is the size of the gap $\mathbf
F - G$ over element $K$ (\cref{eq:zz-indicator}). It is small where the flux is
nearly continuous (left, smooth region) and large where the flux jumps sharply
(centre, where the solution bends most). The estimator never sees the true error---it
infers it from the roughness the discretization left behind.}
\label{fig:zz}
\end{figure}

\subsection{The per-element indicator}

With a smooth flux $G$ in hand, the error indicator on element $K$ is the $L^2$ norm
of the gap, restricted to that element:
\begin{equation}
\boxed{\;
\eta_K^2 = \int_K \norm{\mat F - G}^2 \dV
\;}
\label{eq:zz-indicator}
\end{equation}
We work with the \emph{squared} indicator $\eta_K^2$ throughout, for a reason that
will matter in marking: squared error contributions \emph{add}, and the total squared
error $\sum_K \eta_K^2$ is the quantity the marking step apportions. The square root
that recovers $\eta_K$ itself (with a physical energy normalisation) is a thin
epilogue applied later, when the indicators are compared against a tolerance.

Three properties make \cref{eq:zz-indicator} a good indicator, and each is worth
naming because each is a fact you can check.

\emph{Local and cheap.} The integral in \cref{eq:zz-indicator} is over a single
element $K$, so $\eta_K^2$ is computed element-by-element from the local values of
$\mat F$ and $G$---one small quadrature per element, no global coupling in this step.
(The recovery \cref{eq:zz-recover} of $G$ is a global solve, but it is one solve,
shared by all elements; the per-element \emph{difference} is local.) The cost is a
small constant per element, negligible beside the solve that produced $u_h$.

\emph{Quadratic in the field.} The indicator is a squared norm of $\mat F - G$, and
both $\mat F$ and $G$ are \emph{linear} in $u_h$ (the flux map and the projection are
linear operators). So $\eta_K^2$ is a \emph{quadratic} functional of the solution: it
scales as $\eta_K^2(\alpha u_h) = \alpha^2\,\eta_K^2(u_h)$ but is \emph{not} additive,
$\eta_K^2(u_h + w_h) \ne \eta_K^2(u_h) + \eta_K^2(w_h)$. It is an energy-like
quantity, not a linear measurement.

\emph{Vanishes exactly when the flux is already smooth.} If the computed flux $\mat
F$ already lies in the smooth recovery space---no jumps to recover away---then its
$L^2$ projection $G$ equals $\mat F$ itself, and $\eta_K^2 = 0$ for every element.
The estimator fires precisely on the \emph{non-smoothness} the discretization
introduced, which is exactly the ZZ premise: a big jump means a big error, a smooth
flux means a converged solution.

\begin{pitfall}
The ZZ estimate is an \emph{indicator}, not the error. It is not a rigorous upper
bound on $\norm{u - u_h}$---it can over- or under-estimate the true error by a
problem-dependent constant. What it does reliably is rank the elements: it is large
where the error is large and small where it is small, so it tells you \emph{which}
elements to refine even when it does not tell you the error's exact size. For
\emph{driving refinement} that ranking is all we need (the marking step uses only the
relative sizes); for \emph{certifying} a final accuracy one wants a guaranteed bound,
which is a separate and harder theory. Do not read $\eta_K$ as ``the error on element
$K$''; read it as ``element $K$'s share of the work still to do.''
\end{pitfall}

\subsection{Why recover rather than just measure the jumps?}

One might ask: if the jumps in $\mat F$ are the signal, why project at all---why not
measure the jumps directly, face by face? Both strategies exist in the literature
(the direct one is the ``residual'' or ``jump'' estimator). The recovery approach has
two practical virtues that make it Palace's choice. First, it produces a smooth flux
$G$ that is itself a \emph{better} approximation to the true flux than $\mat F$
was---a useful by-product, sometimes used to improve the reported solution.\ Second,
the projection \cref{eq:zz-recover} is a single mass-matrix solve in machinery the
code already has (\cref{ch:cg}'s conjugate gradients on an SPD mass matrix), so the
estimator reuses the solver infrastructure rather than introducing face-by-face jump
bookkeeping. The estimator is, structurally, ``one extra solve and one cheap
per-element norm''---a small tax on each adaptive pass.

\section{Mark: D\"orfler bulk marking}
\label{sec:mark}

The estimate step hands us a list of nonnegative numbers, one per element:
$\{\eta_K^2\}_{K=1}^{N}$. The mark step must decide \emph{which} elements to refine.
This is a more subtle question than it looks, and the answer---D\"orfler's bulk
marking---is the one choice in the loop with a convergence theorem attached.

\subsection{Three marking strategies, one good one}

The naive strategies are tempting and wrong.

\emph{Refine the top $k$.} Mark the $k$ elements with the largest indicators, for a
fixed $k$. But $k$ is a magic number unrelated to the error: on a mesh where the
error is concentrated in two elements, marking the top $50$ wastes effort; on a mesh
where it is spread over hundreds, marking the top $50$ refines too little.

\emph{Refine the top fraction of elements.} Mark the worst-ranked $10\%$ of
\emph{elements}. This couples to the mesh size but still not to the error: $10\%$ of a
mesh might carry $90\%$ of the error or $5\%$ of it, depending on how concentrated the
features are. Marking by element \emph{count} ignores the very thing we are trying to
equidistribute.

\emph{Refine the elements carrying a fixed fraction of the error.} This is
D\"orfler's idea, and it is the right one. Mark the \emph{smallest} set of elements
whose summed squared error reaches a target fraction $\theta$ of the \emph{total}
squared error:
\begin{equation}
\boxed{\;
\text{find the smallest } \mathcal{S} \text{ with } \;
\sum_{K \in \mathcal{S}} \eta_K^2 \;\ge\; \theta \sum_{K} \eta_K^2.
\;}
\label{eq:dorfler}
\end{equation}
The marking is by \emph{error mass}, not element count. If the error is concentrated
in a few elements, $\mathcal{S}$ is tiny; if it is spread out, $\mathcal{S}$ is large.
Either way the marked set captures the same fraction $\theta$ of the work. The
default fraction in Palace is $\theta = 0.7$: mark the smallest set of elements
responsible for $70\%$ of the total squared error. This is called \emph{bulk}
marking because it marks the bulk of the error, however few or many elements that
takes.

\subsection{The algorithm: sort, accumulate, pivot}

\Cref{eq:dorfler} reads as a combinatorial optimisation---``the smallest set
with''---but it has a one-pass realisation, because greedily taking the largest
indicators first is optimal. The algorithm is:

\begin{enumerate}[nosep]
\item \textbf{Sort} the squared indicators in descending order.
\item \textbf{Accumulate} a running sum down the sorted list.
\item \textbf{Stop} (the \emph{pivot}) at the first point where the running sum
reaches $\theta$ of the total. Everything down to and including the pivot is marked.
\end{enumerate}

That ``largest first'' greedy is optimal for \cref{eq:dorfler}: to reach a fixed
error budget $\theta\sum_K\eta_K^2$ with the fewest elements, you must take the
biggest contributors first---swapping any marked element for a smaller unmarked one
could only require \emph{more} elements to reach the same budget. So the prefix of the
sorted list \emph{is} the smallest set. \Cref{alg:dorfler} writes it in the
framework's pseudo-language; the implementation realises the running sum as a
\emph{prefix sum} and the pivot as a binary search (a \code{lower\_bound}) for speed,
but the meaning is exactly the three steps above.

\begin{lstlisting}[style=l4style,caption={D\"orfler bulk marking: the smallest set covering fraction $\theta$ of the squared error.},label={alg:dorfler}]
dorfler_mark :: Float -> [Float] -> IndexSet
dorfler_mark theta etas2 =
  let total   = sum etas2                       -- total squared error
      ranked  = sortDescendingByIndicator etas2 -- (index, eta2) pairs, biggest first
      target  = theta * total                   -- the error budget to cover
      go acc marked ((i, e2):rest)
        | acc >= target = marked                -- pivot reached: stop
        | otherwise     = go (acc + e2) (i:marked) rest
      go acc marked []  = marked                -- (all elements; only if theta = 1)
  in  go 0.0 [] ranked
\end{lstlisting}

\Cref{fig:dorfler-hist} draws the algorithm as a picture: the squared indicators as
bars sorted tallest-first, and the marked prefix as the leading bars whose cumulative
area first crosses the $\theta$ line. It is the most vivid way to see what bulk
marking does---it walks in from the left, taking the biggest bars, until it has
captured fraction $\theta$ of the total area, then stops.

\begin{figure}[t]
\centering
\begin{tikzpicture}[scale=1.0]
\begin{axis}[
  width=12.5cm, height=6.0cm,
  ybar, bar width=8mm,
  xmin=0.3, xmax=8.7, ymin=0, ymax=46,
  xtick={1,2,3,4,5,6,7,8},
  xticklabels={$K_1$,$K_2$,$K_3$,$K_4$,$K_5$,$K_6$,$K_7$,$K_8$},
  xlabel={elements, sorted by indicator (largest first)},
  ylabel={$\eta_K^2$},
  axis lines=left, tick label style={font=\footnotesize},
  clip=false,
]
  % marked bars (first three: 40+25+14 = 79 >= 70 = theta*total), then unmarked
  \addplot[fill=simBgRust, draw=simRust, thick] coordinates
    {(1,40)(2,25)(3,14)};
  \addplot[fill=simBgGray, draw=simGray] coordinates
    {(4,9)(5,5)(6,3)(7,2.5)(8,1.5)};
  % pivot line between K3 (last marked) and K4 (first unmarked)
  \draw[simGreen, thick, densely dashed] (axis cs:3.5,0) -- (axis cs:3.5,44);
  \node[simGreen, font=\scriptsize, anchor=south] at (axis cs:3.5,44) {pivot};
  \node[simRust, font=\scriptsize\bfseries, align=center] at (axis cs:2,35)
    {\textsc{marked}\\(bulk of the error)};
  \node[simGray, font=\scriptsize, align=center] at (axis cs:6,18)
    {unmarked\\(refine later, if ever)};
\end{axis}
\end{tikzpicture}
\caption[D\"orfler bulk marking on a sorted histogram]{D\"orfler marking, drawn as a
sorted histogram. The per-element squared indicators $\eta_K^2$ are sorted
\emph{descending} and accumulated from the left; the marked set (rust) is the leading
prefix of bars whose cumulative area first reaches fraction $\theta$ of the total. The
pivot (green dashed line) is exactly where the running sum crosses $\theta\sum_K
\eta_K^2$. Here a few tall bars on the left carry most of the error, so a small marked
set suffices---the hallmark of a solution with concentrated features. A flat
histogram (no features) would push the pivot far to the right, marking nearly
everything, which is bulk marking correctly reporting that the error is spread out.
\Cref{wex:dorfler} traces the arithmetic for $\theta = 0.7$.}
\label{fig:dorfler-hist}
\end{figure}

\subsection{The four laws that make marking trustworthy}

The marking rule \cref{eq:dorfler} satisfies four properties that are worth stating
as laws, because the convergence theory of \cref{sec:loop} depends on them and because
they make the rule predictable.

\begin{description}[style=nextline, leftmargin=0pt, font=\normalfont\itshape]
\item[Coverage.] The marked set provably covers at least fraction $\theta$ of the
total squared error: $\sum_{K \in \mathcal S}\eta_K^2 \ge \theta\sum_K \eta_K^2$.
This is the defining inequality, and it is the hypothesis the convergence theorem
needs. The implementation verifies it as a post-condition---it never returns an
under-marking.
\item[Minimality.] Among all sets achieving coverage, $\mathcal S$ has the fewest
elements. This is the ``smallest set'' in \cref{eq:dorfler}, guaranteed by the
largest-first greedy. Minimality is what keeps the mesh from growing faster than it
must: we refine the least we can while still capturing the bulk.
\item[Monotonicity in $\theta$.] A larger target fraction marks a superset:
$\theta_1 \le \theta_2 \implies \mathcal S(\theta_1) \subseteq \mathcal S(\theta_2)$.
At $\theta \to 0^+$ the marked set shrinks to the single largest-indicator element;
at $\theta = 1$ it grows to the whole mesh. The fraction is a smooth dial between
``refine only the very worst'' and ``refine everything (uniform).''
\item[The over-mark tie-break.] The achievable cumulative fractions form a discrete
ladder---one rung per element---so the target $\theta$ usually falls \emph{between}
rungs. When it does, the rule takes the \emph{larger} set: it covers the smallest
fraction $\ge \theta$, never the largest fraction $< \theta$. This ``over-mark rather
than under-mark'' tie-break is load-bearing, not cosmetic---it is what guarantees the
coverage inequality holds with $\ge$, and the convergence theorem assumes exactly
that direction. Flipping it to under-mark would void the guarantee.
\end{description}

\begin{keyidea}
D\"orfler (\emph{bulk}) marking refines the smallest set of elements responsible for
a fixed fraction $\theta$ of the \emph{total squared error}---not a fixed count of
elements, and not a fixed fraction of elements. The algorithm is: sort the indicators
largest-first, accumulate, and stop at the pivot where the running sum reaches
$\theta$. Marking by error \emph{mass} (rather than element count) is what
equidistributes the error, and the ``over-mark on ties'' rule is what makes the
covered fraction provably $\ge \theta$---the single hypothesis behind the convergence
guarantee of \cref{sec:loop}.
\end{keyidea}

\section{Refine: subdividing the marked elements}
\label{sec:refine}

Marking produced a set of element indices. Refinement turns that set into a finer
mesh: each marked element is \emph{subdivided} into smaller children, and the global
finite-element space and operators are rebuilt on the new mesh before the next solve.

\subsection{How an element is split}

Two subdivision strategies dominate, and \cref{fig:refine} draws both for a triangle.

\emph{Red--green refinement.} A ``red'' refinement of a triangle joins the midpoints
of its three edges, splitting it into four congruent children (and a tetrahedron into
eight). This is clean and uniform, but it creates a problem at the boundary of the
refined region: a red-refined triangle has new nodes at its edge midpoints, and if the
neighbouring triangle was \emph{not} refined, those midpoint nodes sit in the middle
of the neighbour's edge with nothing to match them---a \emph{hanging node}. A hanging
node breaks conformity (the mesh is no longer a clean tiling) and must be dealt with.
The ``green'' part of red--green refinement is a closure step: neighbours adjacent to a
red element get a temporary bisection (a single cut to a vertex) that absorbs the
hanging node and restores a conforming mesh, at the cost of some sliver elements that
are undone before the next refinement pass.

\emph{Bisection.} An alternative cuts each marked element in two by a single edge
(the longest edge, or a designated ``refinement edge''), propagating the cut into
neighbours as needed to keep the mesh conforming. Repeated bisection, organised
carefully (newest-vertex bisection), provably keeps element shapes bounded---no cell
degenerates into a sliver however many times it is cut---which is what protects the
conditioning of the assembled matrix.

\begin{figure}[t]
\centering
\begin{tikzpicture}[scale=1.0, every node/.style={font=\footnotesize}]
  % ===== LEFT: red refinement (1 -> 4) =====
  \begin{scope}
    \node[font=\small\bfseries, text=simRust] at (1.0,2.5) {red: $1 \to 4$};
    \coordinate (A) at (0,0); \coordinate (B) at (2,0); \coordinate (C) at (1,1.8);
    \coordinate (mAB) at (1,0); \coordinate (mBC) at (1.5,0.9); \coordinate (mCA) at (0.5,0.9);
    \draw[simInk, thick] (A)--(B)--(C)--cycle;
    \draw[simRust, thick] (mAB)--(mBC)--(mCA)--cycle;
    \fill[simRust] (mAB) circle (1.3pt); \fill[simRust] (mBC) circle (1.3pt);
    \fill[simRust] (mCA) circle (1.3pt);
    \node[simGray, font=\scriptsize] at (1.0,-0.4) {four congruent children};
  \end{scope}
  % ===== MIDDLE: bisection (1 -> 2) =====
  \begin{scope}[xshift=4.0cm]
    \node[font=\small\bfseries, text=simGreen] at (1.0,2.5) {bisect: $1 \to 2$};
    \coordinate (A2) at (0,0); \coordinate (B2) at (2,0); \coordinate (C2) at (1,1.8);
    \coordinate (mid2) at (1,0); % midpoint of longest edge AB
    \draw[simInk, thick] (A2)--(B2)--(C2)--cycle;
    \draw[simGreen, thick] (C2)--(mid2);
    \fill[simGreen] (mid2) circle (1.3pt);
    \node[simGray, font=\scriptsize] at (1.0,-0.4) {cut the longest edge};
  \end{scope}
  % ===== RIGHT: hanging node =====
  \begin{scope}[xshift=8.2cm]
    \node[font=\small\bfseries, text=simBlue] at (1.1,2.5) {hanging node};
    % refined triangle on the left, coarse on the right sharing an edge
    \coordinate (P) at (0,0); \coordinate (Q) at (1.1,0); \coordinate (R) at (0.55,1.7);
    \coordinate (mPR) at (0.275,0.85); \coordinate (mQR) at (0.825,0.85);
    \coordinate (mPQ) at (0.55,0);
    \draw[simInk, thick] (P)--(Q)--(R)--cycle;
    \draw[simRust, thick] (mPQ)--(mQR)--(mPR)--cycle;
    % coarse neighbour to the right sharing edge QR
    \coordinate (S) at (2.0,0.85);
    \draw[simInk, thick] (Q)--(S)--(R);
    % the hanging node sits at mQR, mid of shared edge QR, unmatched in the neighbour
    \fill[simBlue] (mQR) circle (2.0pt);
    \draw[refedge, simBlue] (1.7,1.9) to[bend right=15] (mQR);
    \node[simBlue, font=\scriptsize, align=center] at (1.75,2.05)
      {unmatched\\midpoint};
    \node[simGray, font=\scriptsize, align=center] at (1.0,-0.45)
      {neighbour not refined\\$\Rightarrow$ constrain or close};
  \end{scope}
\end{tikzpicture}
\caption[Refinement of a marked element]{Subdividing a marked triangle. \emph{Left}:
\emph{red} refinement joins the three edge midpoints, producing four congruent
children (eight for a tetrahedron). \emph{Middle}: \emph{bisection} cuts a single
edge (here the longest) into two children, propagating into neighbours to stay
conforming. \emph{Right}: the difficulty either strategy must handle---a refined
element creates a new midpoint node on an edge it shares with an \emph{unrefined}
neighbour, a \emph{hanging node} unmatched on the coarse side. It is resolved either
by a \emph{green} closure (bisecting the neighbour to absorb it) or by a
\emph{constraint} that algebraically ties the hanging degree of freedom to its coarse
neighbours, keeping the global space conforming.}
\label{fig:refine}
\end{figure}

\subsection{Where our framework stops}

The subdivision geometry---which edges to cut, how to propagate cuts, how to generate
and constrain hanging nodes, how many levels of nonconformity to allow---is intricate,
and in Palace it is handled entirely by the underlying mesh library (MFEM). Our
framework treats \code{refine} as an \emph{opaque boundary}: a function
$\code{refine} : \text{Mesh} \to \text{IndexSet} \to \text{Mesh}$ that takes the
current mesh and the marked set and returns a finer conforming mesh, with the
subdivision arithmetic owned by the library. This is a deliberate scope decision, not
an oversight: the estimate and mark stages are short, well-defined algorithms we
specify fully and could re-implement; the refinement kernel is a large, mature piece
of mesh infrastructure we \emph{call} rather than \emph{re-derive}.

\begin{notationbox}
We write the refinement step as \lstinline[style=l4style]|refine mesh marked|,
returning a new mesh. Like every value in this book it is treated as
\emph{immutable}: refinement \emph{produces} a finer mesh rather than mutating the old
one in place (the underlying library does mutate a buffer, but that is below our
boundary, exactly as the in-place arithmetic of a BLAS kernel sits below
\code{linear\_combination}). The new mesh is the value the loop threads forward;
the old mesh is spent.
\end{notationbox}

What matters for the loop is only \code{refine}'s contract: given a marked set, it
returns a mesh that is finer \emph{exactly there}, conforming, and ready for the next
assemble--solve. With that contract we can close the loop.

\section{The loop as a state-generated fold}
\label{sec:loop}

We can now assemble the three stages into the loop of \cref{ch:lifecycle} and see
precisely what kind of computation it is. The shape is a \emph{fold}---the structure
\cref{ch:combinators} made the spine of the whole simulator---but a fold of a special
and important kind: one whose \emph{schedule is generated from its own state}.

\subsection{Recalling the two kinds of fold}

\Cref{ch:combinators} drew the map/fold split: a \code{fold} threads a carry through a
chain of steps, each step consuming what the last produced. The transient analysis of
\cref{ch:time} was the exemplar---it folds a time-step operator over a \emph{fixed,
pre-known} schedule of times $[t_0, t_1, \dots, t_M]$. You know the whole schedule
before you start; the fold simply walks it, threading the field state.

The AMR loop is a fold of the \emph{same combinator}, \code{fold\_solve}, but its
schedule is not fixed in advance. There is no list of meshes laid out before the loop
begins. Instead each iteration \emph{computes the next mesh from the current
solution's error}: solve on the current mesh, estimate the error, mark, refine---and
the refined mesh is the next ``schedule element,'' generated on the spot. The loop
also decides \emph{when to stop} from the state: it halts when the estimated error
falls below tolerance (or a resource budget---a maximum element count or iteration
count---is exhausted). This is a \emph{state-generated} (or \emph{data-dependent})
fold: the schedule is a function of the running carry, not a constant.

\begin{figure}[t]
\centering
\begin{tikzpicture}[font=\footnotesize]
  % ===== LEFT: fixed-schedule fold (time) =====
  \begin{scope}
    \node[font=\small\bfseries, text=simBlue] at (1.8,3.1) {fixed schedule (time-stepping)};
    \node[data, minimum width=9mm] (t0) at (-0.4,2.0) {$s_0$};
    \node[opbox, minimum width=10mm] (k1) at (1.0,2.0) {\code{step}};
    \node[opbox, minimum width=10mm] (k2) at (2.6,2.0) {\code{step}};
    \node[data, minimum width=9mm] (t3) at (4.0,2.0) {$s_M$};
    \draw[flow] (t0) -- (k1);
    \draw[flow] (k1) -- node[above,font=\scriptsize]{$s_1$} (k2);
    \draw[flow] (k2) -- (t3);
    % the pre-known schedule sitting above, feeding the steps
    \node[data, minimum width=30mm, fill=simBgBlue] (sched) at (1.8,0.9)
      {schedule $[t_0,\dots,t_M]$ (known up front)};
    \draw[refedge, simBlue] (1.0,1.55) -- (k1.south);
    \draw[refedge, simBlue] (2.6,1.55) -- (k2.south);
    \draw[flow, simBlue] (sched.north) -- (1.8,1.55);
  \end{scope}
  \draw[simGray!50, thick] (5.3,0.3) -- (5.3,3.4);
  % ===== RIGHT: state-generated fold (AMR) =====
  \begin{scope}[xshift=6.1cm]
    \node[font=\small\bfseries, text=simGreen] at (2.0,3.1) {state-generated (AMR)};
    \node[data, minimum width=11mm] (m0) at (-0.2,2.0) {mesh$_0$};
    \node[opbox, minimum width=13mm, fill=simBgGreen, draw=simGreen] (s1) at (1.6,2.0)
      {\scriptsize solve$\to$est$\to$mark$\to$refine};
    \node[data, minimum width=11mm] (m1) at (3.7,2.0) {mesh$_1$};
    \draw[flow] (m0) -- (s1);
    \draw[flow] (s1) -- (m1);
    % the back-edge: next mesh computed FROM the solve
    \draw[backflow, simGreen] (m1.south) -- ++(0,-9mm) -| node[pos=0.25,below,font=\scriptsize]{if $\eta > \text{tol}$: continue} (s1.south);
    % terminate
    \node[product, minimum width=12mm] (done) at (3.7,0.55) {output};
    \draw[flow] (m1.east) ++(0.0,0) .. controls +(0.6,0) and +(0.6,0) .. (done.east)
      node[midway, right, font=\scriptsize] {$\eta \le \text{tol}$};
    \node[simGreen, font=\scriptsize, align=center] at (1.6,0.95)
      {next mesh is\\\emph{computed}, not listed};
  \end{scope}
\end{tikzpicture}
\caption[Fixed-schedule versus state-generated fold]{Two folds of the same
combinator. \emph{Left}: time-stepping (\cref{ch:time}) folds over a schedule
$[t_0,\dots,t_M]$ that is \emph{known before the loop starts}---the steps read their
times from a pre-built list. \emph{Right}: the AMR loop folds the
solve--estimate--mark--refine step, but its schedule is \emph{generated from the
state}: the next mesh is computed from the current solution's error, and the loop
continues or halts depending on whether the estimated error $\eta$ has fallen below
tolerance. Same carry-threading \code{fold} structure; the difference is whether the
schedule is a constant (left) or a function of the carry (right). This is why the
adaptive loop is a fold with a \emph{data-dependent} stopping condition, not a
\code{map} over a fixed range.}
\label{fig:amrfold}
\end{figure}

\subsection{Writing the fold}

In the framework's vocabulary the loop is one \code{fold\_solve} whose step is the
three-stage body and whose carry bundles the mesh with the running error indicators.
\Cref{alg:amr-fold} writes it. Read it against \cref{fig:amrfold}: the carry is the
mesh (plus the indicators and the scalar error), the step is solve--estimate--mark--
refine, and the loop predicate tests the state.

\begin{lstlisting}[style=l4style,caption={The AMR loop as a state-generated \code{fold\_solve}. The schedule (the sequence of meshes) is generated from the carry, not supplied as a fixed list.},label={alg:amr-fold}]
-- the carry threaded through the fold: a mesh and its error state
type AmrCarry = { mesh : Mesh, etas : [Float], err : Float, ndof : Int }

amr_step :: Estimator -> Float -> AmrCarry -> AmrCarry
amr_step est theta carry =
  let (field, ndof) = solve carry.mesh                 -- assemble + solve (per-driver)
      etas          = flux_recovery_estimate est field -- ESTIMATE (Section 37.2)
      err           = nrm2 etas                         -- global error scalar
      marked        = dorfler_mark theta etas           -- MARK    (Section 37.3)
      mesh'         = refine carry.mesh marked           -- REFINE  (Section 37.4)
  in  { mesh = mesh', etas = etas, err = err, ndof = ndof }

-- the loop: iterate the step WHILE the state says "not yet accurate enough"
adaptive :: Estimator -> RefineConfig -> Mesh -> Product
adaptive est cfg mesh0 =
  let stop carry = carry.err <= cfg.tol            -- data-dependent stop ...
                || exhausted carry.ndof cfg         -- ... or budget hit
      final = iterate_while (not . stop) (amr_step est cfg.theta) (seed mesh0)
  in  product_of final
\end{lstlisting}

Two features of \cref{alg:amr-fold} are the whole point. First, the step \emph{reads
its carry}: \code{amr\_step} solves on \code{carry.mesh}, the mesh the previous step
produced. That read is what makes this a fold and not a \code{map}---the iterations do
not commute and cannot run in parallel or be reordered, because mesh $n+1$ is
\emph{computed from} the solution on mesh $n$. Second, the loop's \code{stop}
predicate is a function of the \emph{state} (the error \code{carry.err} and the
DOF count \code{carry.ndof}), not of a fixed iteration count. The number of
refinement passes is not known in advance; it is however many it takes to drive the
error below tolerance. This is the \code{iterate\_while} of \cref{ch:fold} at its most
characteristic---a loop whose length is discovered, not declared.

\begin{pitfall}
It is tempting to picture the AMR loop as a \code{map} over a list of meshes
``$[\text{mesh}_0, \text{mesh}_1, \dots]$.'' It is not, and the difference is not
pedantic. A \code{map}'s elements are independent and known up front; here neither
holds. Mesh$_{n+1}$ \emph{does not exist} until mesh$_n$ has been solved, estimated,
and refined---it is \emph{produced by} the step, from the carry. You cannot list the
schedule before running the loop, and you cannot run the iterations out of order or in
parallel. The AMR loop is intrinsically a sequential, state-threaded fold; reading it
as a map would promise a parallelism that the data dependence forbids.
\end{pitfall}

\subsection{Why it converges: D\"orfler's theorem}

The loop would be useless if refining did not reliably reduce the error. It does, and
the guarantee is the reason bulk marking, not the naive strategies of
\cref{sec:mark}, is used. The result, due to D\"orfler, is a \emph{contraction}
theorem.

\begin{theorem}[D\"orfler, contraction of adaptive refinement; informal]
\label{thm:dorfler}
For a model elliptic problem, with the error estimator equivalent to the true error
and elements marked by the bulk criterion \cref{eq:dorfler} with fraction $\theta$,
each adaptive pass reduces the error by at least a fixed factor: there is a constant
$\rho = \rho(\theta) < 1$, independent of the mesh, with
\begin{equation}
\norm{u - u_{h}^{(n+1)}} \;\le\; \rho \, \norm{u - u_{h}^{(n)}}
\label{eq:contraction}
\end{equation}
where $u_h^{(n)}$ is the discrete solution after $n$ adaptive passes. Consequently the
error decreases geometrically in the number of passes, and---combined with the
minimality of marking---the adaptive sequence achieves the optimal error-versus-DOF
rate.
\end{theorem}

The content of \cref{thm:dorfler} is that the loop is a genuine \emph{contraction}:
every pass shrinks the error by a guaranteed factor $\rho < 1$, so $n$ passes shrink
it by $\rho^n \to 0$. The bulk-marking coverage law of \cref{sec:mark} is exactly the
hypothesis: because the marked set carries at least fraction $\theta$ of the error,
refining it removes at least a corresponding fixed share, which is what forces
$\rho < 1$. Had we marked by the naive ``top $10\%$ of elements'' rule, no such
guarantee holds---one can construct meshes where the top $10\%$ of elements carry a
vanishing fraction of the error, and refinement stalls. Bulk marking buys the
theorem; the theorem is why the loop is trusted to terminate with a good answer.

\section{Reading the payoff: adaptive versus uniform convergence}
\label{sec:convergence}

The promise of \cref{sec:why-adapt} was that adaptivity reaches a target accuracy with
far fewer unknowns. \Cref{fig:convergence} is where that promise is cashed: a
log--log plot of error against the number of degrees of freedom $N$, for uniform and
adaptive refinement on a corner-singularity problem. It is the single most important
diagnostic an AMR practitioner reads.

On a log--log plot, a convergence rate $\text{error} = C\,N^{-r}$ appears as a
\emph{straight line of slope $-r$}: the steeper the line, the faster the error falls
per unknown added. Two lines tell the story.

\emph{Uniform refinement} (the shallow line) is rate-limited by the singularity. The
corner's $r^{-\gamma}$ blow-up caps the smooth-solution rate: no matter how finely you
cut \emph{every} cell, the under-resolved corner drags the global error down only as
$N^{-\gamma/2}$, a shallow slope. The interior is over-resolved and the corner is
under-resolved, and the corner wins.

\emph{Adaptive refinement} (the steep line) recovers the full rate. By concentrating
cells at the corner---equidistributing the error---it resolves the singularity with a
small, local investment of unknowns and lets the interior stay coarse. The recovered
slope is $N^{-p/2}$, the rate a smooth solution would enjoy. The gap between the lines
is the payoff: read horizontally at any fixed error level, the adaptive mesh needs
\emph{orders of magnitude} fewer unknowns to reach it.

\begin{figure}[t]
\centering
\begin{tikzpicture}[scale=1.0]
\begin{axis}[
  width=12cm, height=7.0cm,
  xmode=log, log basis x=10, ymode=log, log basis y=10,
  xmin=2e2, xmax=3e5, ymin=2e-4, ymax=2e-1,
  xlabel={degrees of freedom $N$}, ylabel={error $\|u-u_h\|$},
  axis lines=left, tick align=outside, grid=major, grid style={simGray!18},
  legend style={at={(0.03,0.03)}, anchor=south west, font=\footnotesize,
    draw=simGray!50},
  every axis plot/.append style={very thick},
  tick label style={font=\footnotesize},
]
  % uniform: slope ~ -1/3 (singularity-limited), shallow
  \addplot[simRust, mark=*, mark size=1.6pt] coordinates {
    (400,1.0e-1)(1600,6.3e-2)(6400,4.0e-2)(25600,2.5e-2)(102400,1.6e-2)};
  \addlegendentry{uniform: slope $-\gamma/2$ (singularity-limited)}
  % adaptive: slope ~ -1 (full rate), steep
  \addplot[simBlue, mark=square*, mark size=1.6pt] coordinates {
    (300,8.0e-2)(900,2.7e-2)(2700,9.0e-3)(8100,3.0e-3)(24300,1.0e-3)(72900,3.3e-4)};
  \addlegendentry{adaptive: slope $-p/2$ (full smooth rate)}
  % horizontal "same accuracy" guide at error = 3.3e-3
  \draw[simGreen, densely dashed, thick] (axis cs:2e2,3.3e-3) -- (axis cs:3e5,3.3e-3);
  \node[simGreen, font=\scriptsize, anchor=south west] at (axis cs:2.2e2,3.5e-3)
    {same accuracy};
  % the DOF gap at that level: adaptive ~7e3, uniform off-chart (extrapolated)
  \node[simGreen, font=\scriptsize, anchor=west] at (axis cs:1.1e4,1.8e-3)
    {$\sim\!10\times$ fewer DoFs};
\end{axis}
\end{tikzpicture}
\caption[Adaptive versus uniform convergence]{The payoff, on a log--log error-vs-DOF
plot for a corner-singularity problem. \emph{Uniform} refinement (rust) is limited by
the singularity to the shallow rate $\|u-u_h\| = \mathcal{O}(N^{-\gamma/2})$---cutting
every cell cannot outpace the under-resolved corner. \emph{Adaptive} refinement (blue)
concentrates unknowns at the corner and recovers the full smooth-solution rate
$\mathcal{O}(N^{-p/2})$, a steeper line. Reading horizontally at any fixed accuracy
(green dashed), the adaptive mesh reaches it with far fewer degrees of freedom---the
gap widens without bound as the target accuracy tightens. This separation, not any
single solve, is what adaptive mesh refinement buys.}
\label{fig:convergence}
\end{figure}

\begin{keyidea}
On a log--log error-versus-DOF plot, a convergence rate is a straight line and a
\emph{steeper line is better}. Uniform refinement on a singular problem is stuck on a
shallow line (the singularity caps the rate); adaptive refinement recovers the steep,
full-rate line by equidistributing the error. The horizontal gap between the lines---
how many fewer unknowns adaptivity needs for the same accuracy---widens as the target
accuracy tightens. That gap is the entire economic case for the
estimate--mark--refine loop.
\end{keyidea}

\section{Worked examples}
\label{sec:worked}

\begin{workedexample}{D\"orfler marking on eight elements, $\theta = 0.7$}{dorfler}
\label{wex:dorfler}
We run the bulk-marking algorithm of \cref{sec:mark} explicitly on a concrete list of
indicators, showing exactly which elements are marked and verifying the coverage law.

\emph{The data.} Eight elements with squared error indicators (in arbitrary units)
\[
\eta^2 = (\,4,\;\;36,\;\;9,\;\;1,\;\;25,\;\;2,\;\;16,\;\;7\,),
\qquad K = 1,\dots,8.
\]
The total squared error is
\[
T = \sum_K \eta_K^2 = 4+36+9+1+25+2+16+7 = 100,
\]
and with $\theta = 0.7$ the error budget to cover is $\theta T = 70$.

\emph{Step 1 --- sort descending.} Rank the elements largest-first, carrying their
original indices:
\[
\begin{array}{c|cccccccc}
\text{rank} & 1 & 2 & 3 & 4 & 5 & 6 & 7 & 8 \\\hline
\text{element } K & 2 & 5 & 7 & 3 & 8 & 1 & 6 & 4 \\
\eta_K^2 & 36 & 25 & 16 & 9 & 7 & 4 & 2 & 1
\end{array}
\]

\emph{Step 2 --- accumulate, and find the pivot.} Walk the sorted list, summing until
the running total reaches the budget $70$:
\[
\begin{array}{c|c|c|c}
\text{take element} & \eta^2 & \text{running sum} & \ge 70? \\\hline
K=2 & 36 & 36 & \text{no} \\
K=5 & 25 & 61 & \text{no} \\
K=7 & 16 & 77 & \textbf{yes} \;\leftarrow\; \text{pivot} \\
\end{array}
\]
The running sum first reaches $70$ when we add the third element, $K = 7$, bringing
the total to $77$. The pivot is here: we stop.

\emph{Step 3 --- the marked set.} The marked elements are the prefix taken up to and
including the pivot:
\[
\mathcal S = \{\,K_2,\;K_5,\;K_7\,\} = \{2, 5, 7\},
\qquad |\mathcal S| = 3.
\]

\emph{Verify the laws.} \textbf{Coverage:} the marked set carries
$36 + 25 + 16 = 77$ of the total $100$, so the covered fraction is $0.77 \ge 0.70 =
\theta$. The coverage law holds. \textbf{Minimality:} could a smaller set cover
$70$? A two-element set tops out at $36 + 25 = 61 < 70$ (the two largest), so no two
elements suffice---three is the minimum, and the greedy found it. \textbf{Over-mark
tie-break:} the achievable cumulative fractions are the ladder
$0.36,\,0.61,\,0.77,\,\dots$; the target $0.70$ falls \emph{between} the second rung
($0.61$) and the third ($0.77$), so the rule takes the larger set ($0.77 \ge 0.70$),
never the smaller ($0.61 < 0.70$). It over-marks, exactly as the convergence theory
requires.

Three of eight elements---the bulk of the error---are marked for refinement; the other
five stay coarse. Notice the marking ignored element \emph{count} entirely: it took
whatever number of elements was needed to reach $70\%$ of the error mass, which here
happened to be three. On a mesh where the error were more concentrated (say all $100$
in one element) it would mark just one; spread evenly over the eight (each
$\eta_K^2 = 12.5$) it would mark six. Bulk marking adapts the set size to the
error distribution.
\end{workedexample}

\begin{workedexample}{A ZZ flux indicator on a 1-D mesh}{zz}
\label{wex:zz}
We compute the Zienkiewicz--Zhu indicator of \cref{sec:estimate} on a tiny
one-dimensional example, showing it is largest where the solution bends most.

\emph{The setup.} Take the interval $[0,4]$ meshed into four unit elements
$K_1=[0,1],\,K_2=[1,2],\,K_3=[2,3],\,K_4=[3,4]$, with a piecewise-linear solution
$u_h$ whose nodal values are
\[
u_h(0)=0,\quad u_h(1)=1,\quad u_h(2)=2,\quad u_h(3)=4,\quad u_h(4)=8.
\]
The field is the (1-D) gradient $\diffop u_h = u_h'$; take material weight $Q=1$ so
the flux is $\mat F = u_h'$. On each element $u_h$ is linear, so $u_h'$ is the
constant slope:
\[
\mat F|_{K_1}=\frac{1-0}{1}=1,\quad
\mat F|_{K_2}=\frac{2-1}{1}=1,\quad
\mat F|_{K_3}=\frac{4-2}{1}=2,\quad
\mat F|_{K_4}=\frac{8-4}{1}=4.
\]
This is the \emph{discontinuous} computed flux: a constant on each element, jumping
$1\!\to\!1\!\to\!2\!\to\!4$ at the interior nodes. The jumps are $0$ (at $x=1$), $1$
(at $x=2$), $2$ (at $x=3$)---the flux is smooth on the left and increasingly broken
on the right, where the solution bends upward.

\emph{Recover a smooth flux.} A simple nodal recovery (the 1-D analogue of the $L^2$
projection \cref{eq:zz-recover}) assigns each interior node the average of the two
adjacent element fluxes, and each end node the one flux it touches:
\[
G(0)=1,\quad
G(1)=\tfrac{1+1}{2}=1,\quad
G(2)=\tfrac{1+2}{2}=1.5,\quad
G(3)=\tfrac{2+4}{2}=3,\quad
G(4)=4.
\]
The recovered flux $G$ is the continuous, piecewise-linear function through these
nodal values---a smooth stand-in for the true flux.

\emph{The per-element indicator.} On element $K$, the indicator
$\eta_K^2 = \int_K (\mat F - G)^2\,dx$ measures the gap between the constant element
flux $\mat F|_K$ and the linear recovered flux $G$. On a unit element from node values
$G_{\text{left}}, G_{\text{right}}$ against the constant $\mat F$, the integral of the
squared difference of a constant and a line evaluates to
\[
\eta_K^2 = \int_0^1\!\Bigl(\mat F - \bigl[G_{\text{left}} + (G_{\text{right}}-G_{\text{left}})\,s\bigr]\Bigr)^2 ds .
\]
Computing each:
\[
\begin{array}{c|c|c|c}
\text{element} & \mat F|_K & (G_{\text{left}},G_{\text{right}}) & \eta_K^2 \\\hline
K_1 & 1 & (1,\,1)   & 0 \\
K_2 & 1 & (1,\,1.5) & \int_0^1(1-(1+0.5s))^2\,ds = \int_0^1 (0.5s)^2 ds = \tfrac{1}{12}\approx 0.083 \\
K_3 & 2 & (1.5,\,3) & \int_0^1(2-(1.5+1.5s))^2 ds = \int_0^1 (0.5-1.5s)^2 ds = 0.25 \\
K_4 & 4 & (3,\,4)   & \int_0^1(4-(3+s))^2 ds = \int_0^1 (1-s)^2 ds = \tfrac{1}{3}\approx 0.333 \\
\end{array}
\]

\emph{Read the result.} The indicators are
\[
\eta^2 = (\,0,\;\;0.083,\;\;0.25,\;\;0.333\,),
\qquad K_1,\dots,K_4,
\]
\emph{increasing left to right}---exactly tracking where the solution bends most. On
$K_1$, where the flux is perfectly constant ($1\!\to\!1$, no jump), the indicator is
zero: the discretization already resolves this region, and the estimator correctly
reports nothing to do. On $K_4$, where the flux jumps hardest ($2\!\to\!4$ across the
node and the solution curves upward fastest), the indicator is largest. Feed these
four numbers to the D\"orfler marker of \cref{wex:dorfler}: with $\theta = 0.7$ and
total $0.666$, the budget is $0.466$; the sorted prefix $0.333 + 0.25 = 0.583 \ge
0.466$ marks $\{K_4, K_3\}$---the two right-hand elements, precisely the region that
needs refining. The estimator found the feature, and the marker selected it, with no
knowledge of the true solution at all.
\end{workedexample}

\section{Where this leaves us}

We have opened the three black boxes of the adaptive loop. \emph{Estimate} is
Zienkiewicz--Zhu flux recovery: smooth the discontinuous computed flux by an $L^2$
projection and take the per-element gap as the indicator---local, cheap, quadratic in
the field, and zero exactly when the flux is already smooth. \emph{Mark} is
D\"orfler bulk marking: sort the indicators, accumulate, and take the smallest prefix
carrying fraction $\theta$ of the total squared error---marking by error mass, with an
over-mark tie-break that makes the coverage provable. \emph{Refine} subdivides the
marked elements (red--green or bisection, keeping the mesh conforming) and rebuilds the
space---a mesh-library operation we treat as a boundary. And the whole loop is a
\code{fold\_solve} whose schedule is \emph{generated from its own state}: the next
mesh is computed from the current solution's error, and the fold runs until the error
falls below tolerance. D\"orfler's contraction theorem guarantees each pass shrinks
the error by a fixed factor, and the log--log convergence plot shows the payoff: the
full smooth-solution rate on a singular problem, at orders of magnitude fewer
unknowns than uniform refinement.

This closes the adaptive loop that \cref{ch:lifecycle} sketched and
\cref{ch:combinators} typed as a state-threaded fold. The error indicators the loop
consumes are themselves \emph{reductions}---a per-element norm of the flux gap, and a
global norm of the indicator vector for the stopping test. Those reductions, and the
broader family of fold-to-a-scalar operations the simulator runs, are the subject of
\cref{ch:reductions}, which we take up next.

\summarysection
Adaptive mesh refinement reaches a target accuracy with far fewer unknowns than
uniform refinement by \emph{equidistributing the error}: it puts small cells only
where the solution's features (corners, interfaces, field concentrations) make the
local error large, leaving smooth regions coarse. It does this automatically through
the estimate--mark--refine loop of \cref{ch:lifecycle}. \emph{Estimate} uses the
Zienkiewicz--Zhu flux-recovery indicator $\eta_K^2 = \int_K \norm{\mat F - G}^2$,
where $\mat F = Q\,\diffop u_h$ is the discontinuous computed flux and $G = \mat
M^{-1}(\text{Flux}\cdot u_h)$ is its smooth $L^2$ projection---a local, cheap,
quadratic indicator that vanishes exactly when the flux is already smooth and is large
where the solution bends most. \emph{Mark} uses D\"orfler bulk marking: the smallest
set of elements whose summed squared error covers a fraction $\theta$ (default $0.7$)
of the total, found by sorting the indicators largest-first and accumulating to a
pivot, with an over-mark-on-ties rule that makes coverage provably $\ge \theta$.
\emph{Refine} subdivides the marked elements (red--green or bisection, kept conforming)
and rebuilds the space---a mesh-library operation at our framework's boundary. The
whole loop is a \code{fold\_solve} whose schedule is generated from its own state: the
next mesh is computed from the current solution's error, and the loop halts when the
error drops below tolerance or a budget is hit---a state-generated fold, not a map
over a fixed schedule. D\"orfler's theorem guarantees geometric error contraction per
pass, and the log--log error-vs-DOF plot shows the payoff: the full smooth-solution
convergence rate on a singular problem, at a fraction of the unknowns uniform
refinement would demand.

\furtherreading
The bulk-marking strategy and its convergence theory are due to
D\"orfler~\citep{dorfler1996}, whose ``convergent adaptive algorithm for Poisson's
equation'' introduced the criterion \cref{eq:dorfler} and proved the contraction
\cref{eq:contraction}; it remains the standard reference and is well worth reading for
the elegance of the marking argument. The flux-recovery error estimator is due to
Zienkiewicz and Zhu~\citep{zienkiewiczzhu1987}, whose ``simple error estimator and
adaptive procedure'' is one of the most cited papers in computational engineering;
the recovered-gradient idea of \cref{sec:estimate} is theirs. For the mathematical
theory tying it all together---a-posteriori estimates, the role of bulk marking in
optimal convergence rates, and the approximation theory behind the rate recovery of
\cref{sec:convergence}---Brenner and Scott~\citep{brennerscott2008} is the standard
graduate reference, building directly on the C\'ea-lemma foundation of
\cref{ch:galerkin}. The mesh-refinement geometry of \cref{sec:refine}---red--green
refinement, newest-vertex bisection, and the hanging-node constraints---is treated in
the finite-element-implementation literature; readers wanting the algorithmic detail
of the boundary we declined to cross will find it there.

\exercisessection
\begin{enumerate}[nosep]
  \item \textbf{Marking by hand.} Six elements have squared indicators
  $\eta^2 = (10,\,2,\,30,\,5,\,18,\,35)$. With $\theta = 0.6$, run the D\"orfler
  algorithm: sort descending, accumulate, find the pivot, and write the marked set.
  Verify the coverage law and confirm minimality (no smaller set reaches the budget).

  \item \textbf{The dial $\theta$.} For the data of \cref{wex:dorfler}
  ($\eta^2 = (4,36,9,1,25,2,16,7)$, total $100$), find the marked set for
  $\theta = 0.3$, $\theta = 0.7$, and $\theta = 0.95$. Confirm the monotonicity law
  $\mathcal S(0.3) \subseteq \mathcal S(0.7) \subseteq \mathcal S(0.95)$. What is the
  marked set as $\theta \to 0^+$? As $\theta \to 1$?

  \item \textbf{Why not ``top $10\%$ of elements''?} Construct a list of $20$
  squared indicators in which the top $10\%$ of \emph{elements} (the two largest)
  carry only a tiny fraction of the total error. Explain, using D\"orfler's theorem
  \cref{thm:dorfler}, why refining those two would stall convergence, and how bulk
  marking avoids the trap. (Hint: make the error spread nearly uniformly.)

  \item \textbf{The ZZ indicator vanishes on smooth flux.} Suppose the computed flux
  $\mat F$ is already continuous and piecewise linear---so its $L^2$ projection $G$
  onto the smooth space equals $\mat F$ itself. Show every $\eta_K^2 = 0$. Then argue,
  conversely, that a single jump in $\mat F$ forces at least one $\eta_K^2 > 0$. Which
  property of the indicator (\cref{sec:estimate}) is this?

  \item \textbf{Quadratic, not linear.} Using $\eta_K^2 = \int_K \norm{\mat F - G}^2$
  and the linearity of the flux map and projection in $u_h$, show
  $\eta_K^2(\alpha u_h) = \alpha^2\,\eta_K^2(u_h)$ but give a concrete two-element
  example where $\eta_K^2(u_h + w_h) \ne \eta_K^2(u_h) + \eta_K^2(w_h)$. Why does this
  mean the estimator is \emph{not} an \code{apply\_linop}-style linear operator?

  \item \textbf{Map or fold?} Explain why the AMR loop of \cref{alg:amr-fold} is a
  \code{fold\_solve} and not a \code{map} over a list of meshes. Identify the carry,
  the step, and the data-dependent stopping condition. Which line of the pseudocode is
  the one that forbids reordering the iterations, and why?

  \item \textbf{Reading the convergence curve.} On the log--log plot of
  \cref{fig:convergence}, the uniform line has slope $-\tfrac13$ and the adaptive line
  slope $-1$. (a)~If a target error of $10^{-3}$ requires $N_{\text{ad}} = 8\times10^3$
  adaptive unknowns, estimate how many \emph{uniform} unknowns the shallow line would
  need to reach the same error by extrapolation. (b)~By what factor does the DOF
  advantage of adaptivity grow each time you tighten the target error by $10\times$?

  \item \textbf{Composite indicators.} A driven (frequency-domain) analysis computes
  \emph{two} ZZ indicators per element---an electric-flux one $\eta_{K,E}^2$ and a
  magnetic-flux one $\eta_{K,B}^2$---and combines them as
  $\eta_K^2 = \eta_{K,E}^2 + \eta_{K,B}^2$. Why is the combination done in the
  \emph{squared} domain rather than on the indicators $\eta_K$ themselves? (Hint: which
  quantity is it that the total error $\sum_K \eta_K^2$ must equal, and how does that
  decompose?) Relate this to the channel-additivity of \cref{ch:reductions}'s
  reduction verbs.
\end{enumerate}
